\label{chap:83-21}

\input{ampliaciones_sucesoras_20260824/parte_i_ii/owners/04ac_geometria_proyectiva_bloch_rev2}

\section{Construction of Hilbert space from the nonadic connection and its phase operators}

\subsection{Matrix realization of each depth}

For \(d\ge1\), define
\[
  \mathcal C_9^{(d)}=(\mathbb Z/9\mathbb Z)^d,
  \qquad
  H_d=\ell^2(\mathcal C_9^{(d)}),
  \qquad
  A_d=B(H_d)\cong M_{9^d}(\mathbb C).
\]
The canonical basis \(\{\delta_x:x\in\mathcal C_9^{(d)}\}\) preserves the
discrete label of each state. The uniform measure induces the normalized
trace
\[
  \tau_d(a)=\frac{1}{9^d}\Tr(a).
\]

The factorization
\[
  \mathcal C_9^{(d+1)}
  =
  \mathcal C_9^{(d)}\times\mathbb Z/9\mathbb Z
\]
produce a unitary identification
\[
  H_{d+1}\cong H_d\otimes\mathbb C^9.
\]
This is the exact input to the operator tower. This does not identify
\((\mathbb Z/9)^3\) with \(\mathbb F_3^6\) as groups: they share \(729\)
elements, but their operations and carry cocycles are distinct.

\subsection{Weyl–Heisenberg Representation of Phase and Displacement on the Triadic Fiber}

Let \(\omega=e^{2\pi i/9}\). On
\(H_1=\ell^2(\mathbb Z/9\mathbb Z)\), we define
\[
  (Xf)(r)=f(r-1),
  \qquad
  (Zf)(r)=\omega^r f(r).
\]
Then
\[
  X^9=Z^9=I,
  \qquad
  ZX=\omega XZ.
\]
The group generated by \(X,Z,\omega I\) is the finite Heisenberg group of
order \(9^3=729\).

\begin{theorem}[Stone--von Neumann finite]
Every irreducible unitary representation of the finite Heisenberg group in
which the center acts by means of the primitive character
\(\omega I\) is unitarily equivalent to the preceding representation.
\end{theorem}

\begin{proof}
Let \(v\) be an eigenvector of \(Z\). The vectors
\[
  v,Xv,\ldots,X^8v
\]
have eigenvalues distinct from \(Z\), since
\(ZX^kv=\omega^kX^kZv\). They are orthogonal and span an invariant
subspace. By irreducibility, they form a basis of the entire representation.
In that basis, \(X\) is the cyclic shift and \(Z\) the diagonal phase.
\end{proof}

\begin{controlbox}
Uniqueness fails if the primitive central character is not fixed or if
reducible representations are allowed. This finite theorem must not be confused
with the continuous version for the canonical commutation relations.
\end{controlbox}

\subsection{Aritmetic Projectors and Incompatibility}

The additive and multiplicative sheets of \(\APP\) define information
subspaces and their orthogonal projectors
\(P_{\Sigma,d},P_{\Pi,d}\). The constructed identities satisfy
\[
 \|[P_{\Sigma,2},P_{\Pi,2}]\|
 =\frac{\sqrt2}{3},
 \qquad
 \|[P_{\Sigma,3},P_{\Pi,3}]\|
 =\frac{\sqrt3}{4}.
\]
Noncommutativity is not deduced from a quantum analogy: it is a finite
matrix theorem. Its "pre-Heisenberg" reading is subsequent to that calculation.

\begin{proposition}
The preceding projectors do not jointly belong to the same
Boolean subalgebra of projections.
\end{proposition}

\begin{proof}
In an abelian subalgebra all projectors commute. The displayed commutators
are nonzero.
\end{proof}

\subsection{Hilbertian structure of the matrix publication and preserved genealogical memory}

Hilbert space provides linear superposition, orthogonal projection,
and spectral calculus. HMT supplies the genealogical index of the basis and distinguishes
which operator represents a transition, a measurement, or a shadow. The
realization does not replace \(\TPK\): it represents it.

\begin{sourcebox}
The APP projectors and the Weyl--Heisenberg pair are part of the preceding
construction. The organization \(H_d,A_d\) and the uniqueness proof are established
here by applying the corresponding classical theorem to the HMT
objects.
\end{sourcebox}


\section{The \(9\) phases of the TPK and continuous geometry}

\subsection{Typed separation between the UHF lift of nonadic ranks and the geometric prolongation of the marked branch}

Let \(p_j=|e_j\rangle\langle e_j|\) be the projector onto a marked gate.
The factorization \(H_{d+1}\cong H_d\otimes\mathbb C^9\) admits two
lifts:
\[
  \iota_d(a)=a\otimes I_9,
  \qquad
  j_{d,j}(a)=a\otimes p_j.
\]

\begin{proposition}[Unit and trace]
For all \(a\in A_d\),
\[
  \tau_{d+1}(\iota_d(a))=\tau_d(a),
  \qquad
  \iota_d(I)=I,
\]
while
\[
  \tau_{d+1}(j_{d,j}(a))=\frac1{9}\tau_d(a),
  \qquad
  j_{d,j}(I)\ne I.
\]
\end{proposition}

\begin{proof}
Since \(\Tr(I_9)=9\) and
\(\Tr(p_j)=1\),
\[
 \tau_{d+1}(a\otimes I_9)
 =\frac{\Tr(a)\,9}{9^{d+1}}
 =\tau_d(a),
\]
and
\[
 \tau_{d+1}(a\otimes p_j)
 =\frac{\Tr(a)}{9^{d+1}}
 =\frac{\tau_d(a)}9.
\]
The identity identities are immediate.
\end{proof}

\begin{intuitionbox}
“Looking at the nine phases of the TPK” means preserving the identity operator on the
complete fiber. An isolated gate preserves a marked history; the nine
together also preserve the unit and the global normalization.
\end{intuitionbox}

\subsection{The UHF limit}

\begin{theorem}[UHF nonadic closure]
The unital inductive limit
\[
  A_{9^\infty}
  :=
  \varinjlim(A_d,\iota_d)
  \cong
  \bigotimes_{n=1}^{\infty}M_9(\mathbb C)
\]
is the UHF algebra of supernatural type \(9^\infty=3^\infty\). It has a
product trace \(\tau\) compatible with all the \(\tau_d\).
\end{theorem}

\begin{proof}
The inclusions are injective, unital, and trace-preserving. The algebraic
union consists of nested finite matrices; its norm closure is,
by definition, a UHF algebra. The product trace is the only trace compatible with the
matrix traces of all the factors \citep{glimm1960}.
\end{proof}

\begin{theorem}[Hyperfinite factor]
Let \((\pi_\tau,\mathcal H_\tau,\Omega_\tau)\) be the GNS representation of the
trace. Then
\[
  \pi_\tau(A_{9^\infty})''\cong\Rhyper,
\]
the hyperfinite factor of type \(II_1\).
\end{theorem}

\begin{proof}
The closure is finite because \(\tau(I)=1\). Extremality of the trace makes it a
factor. It is hyperfinite because the increasing union of the matrices
\(\pi_\tau(A_d)\) is weakly dense. Uniqueness of the separable
hyperfinite factor \(II_1\) completes the identification
\citep{murrayvonneumann1936,connes1976}.
\end{proof}

\subsection{The marked branch changes type}

Let \(V_{d,j}:H_d\to H_{d+1}\), \(V_{d,j}\xi=\xi\otimes e_j\).
Then
\[
  j_{d,j}(a)=V_{d,j}aV_{d,j}^*.
\]
In the limit of the spaces, the union of the corners contains the
finite-rank operators. Its norm closure is
\(\mathcal K(H_\infty)\), and its natural bicommutant is
\(B(H_\infty)\), of type \(I_\infty\).

\begin{figure}[H]
\centering
\resizebox{\textwidth}{!}{%
\begin{tikzpicture}[>=Latex,node distance=12mm and 17mm,
  box/.style={rounded corners,draw=hmtblue,fill=hmtlightblue,
    minimum width=34mm,minimum height=13mm,align=center},
  gold/.style={rounded corners,draw=hmtgold,fill=hmtlightgold,
    minimum width=34mm,minimum height=13mm,align=center}]
  \node[box] (a) {\(M_{9^d}\)\\\scriptsize finite state};
  \node[box,above right=of a] (full) {\(a\otimes I_9\)\\\scriptsize nine phases of the TPK};
  \node[gold,below right=of a] (one) {\(a\otimes p_j\)\\\scriptsize gate marcada};
  \node[box,right=of full] (uhf) {\(\mathrm{UHF}(9^\infty)\)\\\scriptsize unital/tracial structure};
  \node[box,right=of uhf] (ii) {\(\Rhyper\)\\\scriptsize type \(II_1\)};
  \node[gold,right=of one] (k) {\(\mathcal K(H_\infty)\)\\\scriptsize compact};
  \node[gold,right=of k] (i) {\(B(H_\infty)\)\\\scriptsize type \(I_\infty\)};
  \draw[->,very thick,hmtblue] (a)--(full);
  \draw[->,very thick,hmtblue] (full)--(uhf);
  \draw[->,very thick,hmtblue] (uhf)--(ii);
  \draw[->,very thick,hmtgold] (a)--(one);
  \draw[->,very thick,hmtgold] (one)--(k);
  \draw[->,very thick,hmtgold] (k)--(i);
  \node[above=2mm of full,text=hmtgreen]{\scriptsize unit and trace};
  \node[below=2mm of one,text=hmtred]{\scriptsize trace \(\div9\)};
\end{tikzpicture}
}
\caption{A marked door and the nine phases of the TPK produce closures of distinct types.}
\label{p12-83-c21-s2:fig:operator-gates}
\end{figure}

\subsection{Continuous dimension from nonadic ranges}

For a projection \(P\in M_{9^d}\),
\[
  d_d(P)=\tau_d(P)=\frac{\rank(P)}{9^d}.
\]
The values belong to
\(\mathbb Z[1/9]\cap[0,1]\), a dense set in \([0,1]\).

\begin{theorem}[Discrete Continuous Geometry]
For each \(t\in[0,1]\) there exists to projection \(P_t\in\Rhyper\) with
\(\tau(P_t)=t\). Moreover,
\[
  P\sim Q
  \quad\Longleftrightarrow\quad
  \tau(P)=\tau(Q),
\]
where \(\sim\) denotes Murray--von Neumann equivalence.
\end{theorem}

\begin{proof}
Let \(m_d=\lfloor9^dt\rfloor\). Since
\(0\le m_{d+1}-9m_d\le8\), nested projections of rank
\(m_d\) can be chosen. Their strong limit \(P_t\) satisfies, by normality of the trace,
\[
  \tau(P_t)=\lim_d\frac{m_d}{9^d}=t.
\]
The classification of projections in a finite factor gives the final equivalence.
\end{proof}

Continuous dimension is not fractional cardinality. It measures equivalence of
subspaces in the closure. Each approximant preserves integer rank; the
continuity appears when the entire unital genealogy is closed.

\begin{resultbox}
Holographic conservation of the unit becomes a falsifiable
operator-theoretic criterion: replacing \(I_9\) by \(p_j\) divides the trace and changes
the type of the limit. The set of gates is not a census; it is the
condition for unital, trace-preserving inclusion.
\end{resultbox}


\section{Profinite dynamical system: translation, logical algebra, and state space}

\subsection{The 2nd route to the hyperfinite factor}

The HMT construction produces the compact group
\[
  K_\infty\cong C_4\times\mathbb Z_3
\]
as a limit of Pell norm kernels, together with a minimal and uniquely
ergodic translation \(T_u(x)=x+u\) for the Haar measure
\(\mu_H\). We consider
\[
  M_{\rm clk}
  =
  L^\infty(K_\infty,\mu_H)\rtimes_{T_u}\mathbb Z.
\]

\begin{theorem}[Clock factor]
If \(T_u\) is free, ergodic, and Haar-preserving, then
\[
  M_{\rm clk}\cong\Rhyper.
\]
\end{theorem}

\begin{proof}
The free and ergodic p.m.p. action produces a finite trace factor.
\[
 \tau_{\rm clk}\!\left(\sum_nf_nV^n\right)
 =\int f_0\,d\mu_H.
\]
Amenability of \(\mathbb Z\) makes the orbit relation hyperfinite;
therefore, the factor is the hyperfinite \(II_1\).
\end{proof}

\begin{figure}[H]
\centering
\resizebox{\textwidth}{!}{%
\begin{tikzpicture}[>=Latex,node distance=12mm and 15mm,
  box/.style={rounded corners,draw=hmtblue,fill=hmtlightblue,
    minimum width=36mm,minimum height=13mm,align=center},
  boundary/.style={rounded corners,draw=hmtred,dashed,fill=white,
    minimum width=36mm,minimum height=13mm,align=center}]
  \node[box] (tower) {\((\mathbb Z/9)^d\)\\\(M_{9^d}\)};
  \node[box,right=of tower] (uhf) {\(\mathrm{UHF}(9^\infty)\)};
  \node[box,right=of uhf] (r1) {\(\Rhyper\)};
  \node[box,below=18mm of tower] (clock) {\(C_4\times\mathbb Z_3\)\\\(T_u\)};
  \node[box,right=of clock] (cross) {\(L^\infty(K_\infty)\rtimes\mathbb Z\)};
  \node[box,right=of cross] (r2) {\(\Rhyper\)};
  \node[boundary,right=16mm of r1,yshift=-15mm] (theta)
    {\(\Theta\) Canonical HMT\\\scriptsize measure, carry, phase};
  \draw[->,thick,hmtblue] (tower)--(uhf);
  \draw[->,thick,hmtblue] (uhf)--(r1);
  \draw[->,thick,hmtblue] (clock)--(cross);
  \draw[->,thick,hmtblue] (cross)--(r2);
  \draw[<->,very thick,hmtgreen] (r1)--node[right]{\scriptsize isomorfia abstracta}(r2);
  \draw[dashed,->,hmtred] (r1)--(theta);
  \draw[dashed,->,hmtred] (r2)--(theta);
\end{tikzpicture}
}
\caption{Two abstract routes toward \(\Rhyper\). HMT causal conjugation
is an additional problem.}
\label{p12-83-c21-s3:fig:two-routes-r}
\end{figure}

The coincidence of weak closures does not identify the antecedents
\(C^*\). The UHF branch and the odometer crossed product may have
distinct \(K\)-theoretic invariants. A canonical HMT identification requires
a map \(\Theta\) with
\[
  \Theta_*\mu_{\APP}=\mu_H,
  \qquad
  \Theta S_{\rm carry}=T_u\Theta,
\]
and preservation of diagonals, phase, orientation, and projectors. That
intertwiner remains as boundary.

\subsection{theorem ergodico medio of von Neumann for the odometro of Haar in \(K_\infty\)}

Let \(U f=f\circ T_u^{-1}\) in \(L^2(K_\infty,\mu_H)\). The theorem
ergodico medio of von Neumann \citep{vonneumann1932ergodic} gives
\[
 \frac1N\sum_{j=0}^{N-1}U^jf
 \longrightarrow
 P_{\mathrm{Fix}(U)}f.
\]
The ergodicity identifies the limit with
\[
  \left(\int f\,d\mu_H\right)\mathbf1.
\]
The clock has pure point spectrum and zero entropy, but its temporal
averages converge to the Haar average. Mixing is not required.

\subsection{Orthomodular Lattice of Projections and the Logic of Closed Subspaces}

The lattice \(\Proj(\Rhyper)\), with
\[
  P^\perp=I-P,
  \qquad P\wedge Q,
  \qquad P\vee Q=(P^\perp\wedge Q^\perp)^\perp,
\]
is complete and orthomodular. It is not distributive. The
noncommuting APP projectors embed into \(\Rhyper\) and provide a finite
witness of that nondistributivity \citep{birkhoffvonneumann1936}.

Each Abelian subalgebra selects a Boolean context. Classical logic does
not disappear: it occupies the commutative contexts. What fails globally is
the possibility of gathering all the incompatible observables into a single
Boolean algebra.

\subsection{Density matrices and entropy}

In \(M_{9^d}\), to normal state is written
\[
  \varphi_\rho(a)=\Tr(\rho a),
  \qquad \rho\ge0,\quad\Tr\rho=1.
\]
In \(\Rhyper\), a normal density is
\[
  h\in L^1(\Rhyper,\tau)_+,
  \qquad \tau(h)=1,
\]
and \(\varphi_h(a)=\tau(ha)\). Its entropy relative to the trace is
\[
  S_\tau(h)=-\tau(h\log h).
\]

If \(h=9^d\rho\), then
\[
  S_\tau(h)=S_{\rm vN}(\rho)-d\log9.
\]
The canonical extension
\[
  \rho_{d+1}=\rho_d\otimes\frac{I_9}{9}
\]
adds \(\log9\) to the matrix entropy, but leaves
\(S_\tau(h)\) invariant. The first quantity records an additional uniform fiber;
the second uses the holographic normalization of the tower.

\begin{controlbox}
\(\Tr\) and \(\tau\) will not be used interchangeably. Nor will
\(\tau(P)\) be interpreted as a cardinality of points, or every vector in a
representation be called pure. A diffuse factor \(II_1\) has no minimal projections.
\end{controlbox}

\subsection{Operator types and modular boundary}

The classification constructed so far is:
\[
\begin{array}{c|c}
\text{object}&\text{type}\\
\hline
M_{9^d}&I_{9^d}\\
\text{marked branch}&I_\infty\\
\text{tracial UHF branch}&II_1\\
\text{p.m.p. clock}&II_1.
\end{array}
\]
A type \(III\) would require a quasi-invariant measure with nontrivial
Radon--Nikodym cocycle or a KMS state whose GNS representation were not
semifinite. The mere presence of the word “KMS” does not select the type.


\begin{workbox}
The completion of the total memory groupoid of \TPK\ requires a topology
or Borel structure, measure, closure under inversion, amenability,
ergodicity, and control of isotropy. The clock subgroupoid satisfies those
hypotheses; the total groupoid is not promoted by analogy.
\end{workbox}
